%!TEX root = ../thesis.tex
%*******************************************************************************
%****************************** Third Chapter **********************************
%*******************************************************************************
\chapter{Discussion}

% **************************** Define Graphics Path **************************
\ifpdf
    \graphicspath{{Chapter4/Figs/Raster/}{Chapter5/Figs/PDF/}{Chapter5/Figs/}}
\else
    \graphicspath{{Chapter5/Figs/Vector/}{Chapter5/Figs/}}
\fi

A fundamental question motivated this thesis: Can the utility signal encoded by dopamine neurons in the midbrain in nonhuman primates (\cite{staufferDopamineRewardPrediction2014}) be correlated with the economic choices made at the precise moment of decision? 

The opportunities to record from these neurons are rare in humans (\cite{zaghloulHumanSubstantiaNigra2009}). Therefore, researchers face the dual challenge of communicating complex economic task structures to subjects while simultaneously capturing their subjective utility assessments for different goods. This raises a critical methodological question. How can we obtain both neural and behavioural measures of utility concurrently?

\clearpage
\section{The Becker-DeGroot-Marschak auction mechanism allows subjects to reveal their utility for a good at the time of the trial}

In Chapter 2 we outline a well-studied task from the behavioural economics literature: the Becker-DeGroot-Marschak auction mechanism, which offers a compelling alternative to the binary choice tasks which have traditionally been employed in studies of non-human behaviour economics.

In the BDM, subjects bid their willingness-to-pay for goods, with the optimal strategy being to bid exactly one's true utility, eliminating strategic considerations that might otherwise distort truthful reporting. We provide a mathematical treatment of the tasks, showing a proof of the dominant, truthful strategy.

A mathematical analysis of the BDM payoff structure and its implications for animal training revealed a fundamental challenge to the implementation of the BDM task: the shallow cost of misbehaviour function that characterizes second-price auctions. Even substantial deviations from optimal bidding result in relatively small utility losses, creating weak reinforcement signals for learning. For example, a monkey bidding randomly across the entire budget range would forfeit only a small fraction of potential utility compared to optimal performance, making it difficult for animals to distinguish good from poor strategies through experience alone. We further review other potential limitations of the BDM auction task and it's use in human economic decision-making studies (see: \cite{casonMisconceptionsGameForm2014}, \cite{hamukwalaDesignFactorsInfluencing2019}).

Theoretical results are presented, deriving the payoff a behaving monkey could expect was derived given different computer bidding strategies, comparing uniform, normal, and matched normal distributions. The uniform distribution was selected for its analytical tractability and consistent expected payoffs, though normal distributions created sharper cost gradients around optimal bids. Our analysis, combined with data from a previous training study (\cite{almohammadRewardValueRevealed2022} showed that monkeys could expect to receive substantial liquid rewards, over 1ml per trial on average.

The theoretical results in this chapter informed our implementation of the behavioural task that we developed for the monkeys. For any future studies employing auction tasks where the payoff structure is critical to gathering sufficient data (i.e., especially additional animal studies), the results should provide a framework for researchers to calibrate their experiments.

\clearpage
\section{Non-human primates can be trained to perform many sequential auction trials for a juice reward}

In the third chapter of this thesis, we analyse the results of the training of nonhuman primates on the BDM auction task described in Chapter 2. Two rhesus macaque monkeys who had previously been trained on a similar task (\cite{almohammadRewardValueRevealed2022}) were retrained on a version of the BDM optimised for neuronal recording.

Both animals showed a remarkable ability to rank their mean bids according to the three levels of goods offered per session. Although mean bids were ordered in almost all sessions, \textit{individual bids} for the same good offered varied from trial to trial. The cause of this variation is likely due to two distinct causes:

\begin{enumerate}
    \item changing satiety for the juice and water offered per trial over the course of an experimental session
    \item the stochastic nature of economic choices under a random utility model of decision-making.
\end{enumerate}
The ability to separate these two causes of variability is a crucial advantage to the BDM task in comparison to more traditional 'binary choice' experimental paradigms, also described within this thesis. These tasks require choices to be averaged over blocks of trials to back-calculate the value of an offered good, and so are less able to capture these moment-to-moment changes in a subject's utility function.

In this chapter, we also show data from the same two monkeys completing sessions of a binary choice task that is mathematically (but not operationally) identical to the BDM task. In both tasks, non-human primates were able to complete hundreds of trials per session over the course of hours, which would be impractical in human studies.

We compare the behaviour elicited by both the BDM and binary choice tasks and show that the monkeys produced consistent bidding patterns. However, the willingness-to-pay on the BDM sessions was generally slightly higher than in the binary choice task, with one notable exception: at the highest reward level, Monkey V's willingness-to-pay was substantially \textit{greater} on the binary choice task, with an indifference point that could not be directly measured within the tested range. Subjects (neither human nor animal) are not hyper-rational economic agents; the mechanism of extracting a report of utility affects its value.

This general pattern of behaviour (auction WTP exceeding choice WTP) itself runs counter to the usual direction reported in the human auction-versus-choice literature discussed in Chapter 3. Chapter 2's discussion of the endowment effect would in fact predict the more usual, literature-consistent direction: a subject endowed with a water budget should, if anything, be reluctant to spend it, suppressing rather than inflating BDM bids relative to choice. Monkey V's reversal at the highest reward level is the one condition that does follow this more conventional direction, though it should be treated cautiously: the animal's binary choice sessions were curtailed by the early, unexpected failure of its cranial implant (Chapter 3), leaving this particular indifference point estimated from substantially fewer trials and sessions than any other condition.

These bundle-choice-derived willingness-to-pay estimates also let us ask how closely the monkeys' BDM bids approached this ground truth. Using a logistic regression of binary choices to identify reversal values, we model the optimality of primates' BDM choices. Both animals tended to bid within 2\% of the modelled optimal behaviour.

Finally, we present data on how monkeys produce their bids. We analyse the joystick manipulation of the bidding indicator by the monkeys, showing two distinct patterns of bidding: either quickly moving to a desired bid, or dithering with smaller, slower movements towards a final offer.

As expected, the reward offer was by far the dominant variable affecting the monkey's valuation of an offered good in a BDM auction. As expected and indeed part of the motivation for this study, there appeared to also be slight satiety effects, with bids increasing as the session progressed (as monkeys value taste over total thirst-quenching liquid volume). This was confirmed in experimental sessions for both monkeys in which only one reward-representing fractal was offered to control for the magnitude of juice offered.

These single-fractal sessions also reveal the first of several points across this thesis where the two monkeys diverge rather than replicate one another. When only the medium-value fractal was shown throughout a session, Monkey V's bids for it were reliably lower than when the same fractal was shown alongside the high and low fractals (Cohen's $d = 0.467$), while Monkey U showed a negligible difference ($d = 0.077$). This is most naturally read as a relative-valuation effect rather than a magnitude effect, where removing the highest and lowest value fractals from a session changes the distribution of rewards against which the medium fractal's value is judged. This is consistent with range-adaptive value coding described elsewhere in both primate neuroeconomics and behavioural economics more broadly.

To test whether this effect is mirrored neurally, we compared the dopamine response to the medium fractal between single and multi-fractal sessions. Because different, non-overlapping cell populations were necessarily recorded in each session type, this comparison is between rather than within cells, reducing its sensitivity independent of sample size. A mixed-effects model found no significant effect of session context for either monkey, and per-cell Bayes factors favoured the null for both, exceeding the conventional threshold for 'some evidence'. Responses still tracked bid within single fractal sessions for Monkey V, but their context-dependence was not replicated neurally in either animal. This suggests the dopamine signal is robust to, rather than adapted to, session context, at least for the medium fractal.

In comparison to auction sessions, regression analysis of binary choices showed similar dominance of juice and water magnitude variables in determining the monkey's decisions, with minimal influence from trial history or session progression effects. Monkey U showed a negligible bias towards choosing the right-hand option, while Monkey V showed a significant, though still modest, bias towards the left.

\clearpage
\section{Individual dopaminergic neurons behave as if they are encoding economic utility when a non-human primate makes a bid for a desired good}

In our final experiment, we took the BDM paradigm which we had investigated mathematically and behaviourally, and studied the neurobiology of non-human primates making economic auction decisions. Monkeys were implanted with cranial windows and trained to perform the BDM task while head-fixed. The recording location of putative dopaminergic midbrain neurons was confirmed, and stereotypical reward prediction responses elicited by the unpredicted release of small volumes of rewarding water.

Over the course of the study, we successfully recorded from 268 putative dopamine neurons (123 from Monkey U, 145 from Monkey V) that displayed the characteristic electrophysiological properties expected of midbrain dopaminergic cells. These neurons showed the canonical two-component response pattern to reward-predicting stimuli that has been well documented in the dopaminergic literature (\cite{schultzDopamineRewardPredictionerror2016}). The first component reflected sensory detection and attention towards salient reward-associated stimuli, while the second component was associated with valuation processing and showed bidirectional responses dependent on reward magnitude.

Consistent with previous established findings from dopaminergic neurons, a substantial proportion of the cells we recorded demonstrated clear rank-ordered responses to the physical reward magnitude signalled by the three fractals. Population-level analysis revealed highly significant differences in neuronal activity across the three reward levels and at the individual neuron level, 31 of 123 cells in Monkey U and 41 of 145 cells in Monkey V showed significant rank-ordered responses to the discrete reward value offered in a trial. This replicates the well-characterized finding that dopamine neurons encode the magnitude of physical rewards offered to the animal (\cite{toblerAdaptiveCodingReward2005, fiorilloDiscreteCodingReward2003}), providing validation of our recording methodology.

Critically, when dopamine neuron responses were analysed in relation to the monkeys' actual bids rather than simply the objective reward magnitude offered, we observed a striking correspondence between neural activity and subjective valuation. Individual neurons showed changes in firing rate related to the monkey's forthcoming bid during the fractal offer epoch, with 73 cells (32/123 in Monkey U, 41/145 in Monkey V) demonstrating significant correlations between cellular responses and the bid on a linear regression. This relationship held even when controlling for the objective reward magnitude, suggesting that these neurons encode subjective utility rather than simply the sensory properties of the stimulus. This held for five of the six monkey-by-fractal combinations once a mixed-effects model corrected for cell-level clustering; the exception was Monkey V's response to the highest-value fractal, which showed the strongest such relationship under the naive, trial-pooled regression but was no longer significant once this correction was applied.

The most compelling evidence for utility encoding came from trials where monkeys made similar bids for different reward magnitudes. In these cases, dopamine neurons showed remarkably similar activity patterns, indicating they were tracking the monkey's internal valuation of the offered good rather than the external physical magnitude associated with it. We quantified this using Bayes factors comparing firing rates between matched-bid trial pairs drawn from different reward magnitudes: five of the six monkey-comparison combinations showed no evidence for a residual magnitude effect, with the sixth only marginally exceeding the conventional threshold for 'some evidence' -- a genuinely equivocal, rather than decisive, result.

Population-level analysis also revealed significant correlations between neural activity and forthcoming bids (Monkey U: R² = 0.85; Monkey V: R² = 0.88), with relationships remaining when neuronal responses were normalised to baseline activity levels. Because trial-level regressions of this kind pool many trials recorded from the same cell as though they were independent observations, we additionally confirmed this relationship using a mixed-effects model with a random intercept per cell, fitted directly on the continuous bids made per trial; the relationship held for both monkeys (Monkey U: β=0.781±0.070, p=9.70×10⁻²⁹; Monkey V: β=0.644±0.053, p=1.68×10⁻³⁴), confirming that the central bid–dopamine relationship is not an artefact of treating non-independent trials as independent observations.

Importantly, the correlation between neural activity and the bids made occurred during a time window before the monkeys had begun to physically place their bids, suggesting that dopamine neurons encoded the formation of subjective value assessments rather than simply reflecting the execution of the motor response necessary for the offer. Our findings provide direct evidence that dopaminergic cells in the midbrain respond as if they encode the utility of possible economic choices in a precise and instantaneous manner.

\subsection{Individual dopaminergic neurons show stereotyped reward prediction error responses across the BDM task}\label{subsec:stereotyped_rpe}

The BDM auction mechanism provides a rich task which allows us to investigate the behaviour of these recorded putative dopaminergic neurons beyond just their encoding of subjective value, and future bid, at the moment of reward offer. At the trial outcome epoch, when monkeys learned whether they had won or lost the auction, these same neurons displayed the canonical reward prediction error responses that have been extensively documented in the dopaminergic literature.

In contrast to the deterministic binary choice task we presented behavioural results from, the auction task is inherently probabilistic, with trial outcomes dependent on the randomly generated computer bid. After making a bid, this bid, and the outcome of the trial, is revealed to the monkey. This naturally produces a reward prediction error; the monkey can only 'profit' on a trial when they win (and the computer's bid for a good is lower than their own). The neurons we recorded reflect this, with 72 of 123 neurons in Monkey U and 46 of 145 neurons in Monkey V, demonstrating that individual cells were sensitive to the binary outcome of each auction trial.

This encoding extended beyond simple win/lose coding. When monkeys won auctions, the variance in neuronal activity was exceptionally well explained by the actual profit earned on each trial (Monkey U: R² = 0.923; Monkey V: R² = 0.913). Dopamine neurons were responding not just to whether the monkey had won, but also precisely how much economic benefit they had gained from their successful bid. In contrast, when monkeys lost trials, and no profit nor loss is accrued by the animal, a naive trial-level regression found no significant relationship between neuronal activity and the computer's bid for either monkey. This changes, however, once the repeated-measures structure of the data (many trials recorded per cell) is properly accounted for with a mixed-effects model: Monkey V, though not Monkey U, then shows a significant negative relationship between activity and the computer's bid. This is a striking result since the computer's bid has no bearing on a monkey's payout once a trial is lost. We interpret this as a missed-opportunity, or regret-like, signal. A low losing computer bid means the auction could have been won cheaply had the monkey bid a little higher, a possibility we discuss further in Chapter 4 in relation to the anticipated-regret literature and this thesis's own regret framework (Chapter 2).

A similar sensitivity to the computer's bid appears behaviourally, too: Chapter 3 found that monkeys' own subsequent bids were higher following a higher randomly drawn computer bid on the previous trial, and also higher following a win on the previous trial. The latter of these runs counter to the 'hot hand' style prediction that a recent win should instead lead to lower, more confident bidding. Across both chapters, the computer's bid appears to carry more behavioural and neural weight than a strictly incentive-compatible reading of the BDM would predict. The encoding of utility across the trial by the putative dopaminergic neurons recorded was not just a 'one-and-done' representation of the offered good, but it is tracking the subjective valuation by the monkey across the trial.

This tracking continues into the final epochs of the task where the liquid rewards are actually paid out to the monkeys. The receipt of the juice reward led to an increase in the firing rate of putative dopaminergic neurons for both animals, with the greatest changes seen when higher magnitude goods were delivered. Splitting the responses by tercile of the budget, the monkeys placed their bid showed a rank ordering of the change in neuronal activity, suggesting that the monkeys are not simply responding to the volume of rewarding liquid received, but it's subjective value.

A naive regression additionally found that, for Monkey V but not Monkey U, the change in normalised firing rate of putative dopaminergic neurons at the payment of the juice reward was correlated with the preceding bid. This relationship did not survive the same mixed-effects correction described above, however, and is treated as a false positive of the naive, trial-pooled analysis rather than a genuine effect; the tercile-based evidence for subjective value tracking at this epoch, described immediately above, is unaffected by this correction.

Finally, upon receipt of any remaining budget water, the activity of putative dopaminergic neurons was significantly correlated with the total volume of liquid received (and not the 'profit' in terms of the water budget made by the subject on the trial). This utility signalling occurred even though monkey's theoretically had perfect information about the forthcoming water payment from the moment the computer bid is revealed. We discuss why these responses persist despite the canonical extinction of reward prediction errors (\cite{schultzNeuralSubstratePrediction1997}) as an animal learns to expect a reward. Namely, the substantial time delay between the signalling of trial outcome and the heterogeneity of individual BDM trials, which makes such calculations difficult (\textbf{see Section 5.4.3}).

\subsection{Individual differences between subjects}

One further point is worth noting here, a natural consequence of working with a sample of two animals: as noted in Chapter 3, replication across two animals does not imply that the subjects are interchangeable. The monkeys both replicate the central claims of this thesis, but diverge from one another at several points on secondary measures.

The single fractal relative valuation effect discussed above is one example: Monkey V's bids for the medium fractal are lower when it is the only fractal on offer, while Monkey U shows no comparable effect. A second, and perhaps clearer, example comes from what drives bidding within those same single fractal sessions. Monkey U's bids there are dominated by anchoring to the randomly drawn starting bid position, with no significant effect of accumulated liquid; Monkey V's bids are instead dominated by accumulated liquid received, with a comparatively weak starting-bid effect whose sign is reversed relative to Monkey U's.

A neuronal example is less straightforward to interpret. At the reveal of a losing trial, the population of neurons recorded from Monkey U shows a decrease in activity, while Monkey V's shows an increase. Taken at face value this looks like a genuine difference in how the two animals' dopaminergic systems process a loss. However, when the same loss trials are instead split by the tercile of the monkey's own bid, both animals show the same underlying pattern: an increase in activity specific to trials on which a low bid was made. This suggests the population-level `reversal' more likely reflects a difference in how each monkey's losing trials happen to be distributed across bid levels, rather than two distinct computations, and we would caution against too strong a reading of this apparent reversal. None of this bears on the central claims of this thesis that dopaminergic activity is rank-ordered by reward magnitude and correlated with the monkey's own bid, which hold for both animals throughout. This section is included as a reminder that every result reported here was produced by two individual animals, each with their own internal life, not a single averaged subject.

\clearpage
\section{Contribution to the dopaminergic neuroscience literature}

There is no 'hedonimeter' which can tell a researcher the pleasure or pain experienced by an organism for a given outcome. Therefore, previous studies of nonhuman primate decision-making have relied on the calculation of an animal's utility on a given choice trial after the fact and from aggregate choices. However, it seems obvious from observing animals (and from lived human experience) that many choices are made almost on a whim and in the moment; it is not clear that we should think that artificial economic choices made it a laboratory setting are any different.

Here, we present evidence that dopaminergic responses to offered choice can track the subjective value captured by an animal within the economic options present on such a moment-by-moment basis. This is not a surprising finding; what else would we expect a nervous system capable of making instantaneous decisions to be doing? We feel it is, however, an important one. Many of the most important choices made outside of the laboratory are single-shot, and to properly understand these we cannot rely purely on averaged trial responses.

Finally, this and a previous study which trained non-human primates on the BDM task (\cite{almohammadRewardValueRevealed2022}) make a substantial contribution to the behavioural literature. Second price auctions are conceptually difficult for humans to grasp (\cite{casonMisconceptionsGameForm2014}), there was substantial doubt that it was possible to teach them to non-human primates. Our studies provide a proof of concept that it can be done. What other potentially informative tasks have yet to be attempted in non-human models? Recent studies have employed other complex paradigms such as the Prisoner's dilemma (\cite{haroushNeuronalPredictionOpponents2015}, \cite{vancoeverdenFormalCooperationMacaque2022}) to investigate social cognition, and transitive inference tasks which rely on metacognition (\cite{dibelloPrefrontalCortexContribution2024}) in non-human primates\footnote{We touch here on studies in which macaques are the model organism as with our auction. Researchers should be similarly ambitious about using phylogenetically more distant species as well.}. By not limiting the contexts in which we study the nervous systems of model animals, we have the opportunity to reap a much greater understanding of how they work.

\subsection{Limitations of our study and future directions of  research}

Despite our belief that the research presented in this thesis represents a meaningful contribution to the field, no study is without limitations. We identify three opportunities in which our findings could potentially be extended without any particularly complex changes to the either the tasks or training described.

The first potentially interesting avenue of future research that we did not manage to investigate is comparing the activity of single recorded putative dopaminergic neurons between the contexts of our BDM auction task and the mathematically (but not mechanistically) identical binary choice task. In Chapter 3 we present behavioural data comparing the two tasks which both monkeys were highly trained on.

Different reward dimensions, such as probability, payoff time (\cite{lakDopaminePredictionError2014}), and even effort (\cite{burrellWorthWorkMonkeys2023}), are known to be integrated into the dopamine reward signal. Recent work has also shown that dopamine reward prediction errors reflect hidden state inference, with different patterns of activity depending on the structure and uncertainty of the task (\cite{starkweatherDopamineRewardPrediction2017}). Even if the expected payout on a trial is identical between the BDM auction and a binary choice, the effects of needing to consider and bid exactly what one's subjective value for a good is may create different computational demands and state uncertainties that could be encoded differently in dopaminergic activity between the two contexts. For example, dopaminergic cells are known to signal the belief in the accuracy of a perceptual choice (\cite{lakMidbrainDopamineNeurons2017})- would they also signal the belief of the animal that they had accurately reported their subjective value on an auction task? When signalling a binary choice between two outcomes, tasks generally do not differentiate between the physical movements made by the studied animal and so there is no need to be precise beyond any threshold of a valid choice. Meanwhile, in the BDM task the effort required by the animal is proportional to change from the starting bid they wish to make. Within this study we have treated this effort as negligible; neither the gross nor net movement of the bidding indicator, nor the time taken to place a bid, was significantly correlated with the residual of the bid-firing rate regression reported in Chapter 4, suggesting that motor effort does not confound the central bid-dopamine relationship reported here. However, effort was not systematically manipulated in this task, and it would be possible to do so directly, studying whether dopamine neurons integrate these costs into the final value signal or if they independently track the 'precision' of the animal's report of subjective value.

Even without changing the task mechanism, investigating how probability and magnitude variations affect reward processing would be valuable. We kept the rewards offered and their visual cues constant throughout the recording to compare many trials of well-learned reward magnitudes that may vary in subjective value over time. However, the trial-by-trial value reports of the BDM task create an opportunity to study how dopaminergic systems learn and encode novel reward parameters. Dopamine neurons are known to shift their responses from reward delivery to predictive cues during learning (\cite{ljungbergResponsesMonkeyDopamine1992}), and the BDM paradigm would allow simultaneous measurement of this neural learning process alongside the animal's real-time valuation of novel stimuli.

Most simply, but perhaps most challenging, capturing more trials per recorded dopaminergic cell would also greatly benefit future studies. Although monkeys performed an average of 140 auction trials per head-fixed behavioural session (Table \ref{tab:behaviour_trials}), this dropped to approximately 115 BDM trials per recorded cell. This reduction likely resulted from satiety following free rewards and boredom during the time required to locate and confirm recorded cells. In our initial retraining of the monkeys, we were able to frequently elicit over 200 completed trials per day. Spending a greater period of time ensuring that monkeys are well-trained to complete more tasks, and that researchers are able to more quickly find and identify dopaminergic cells when recording, would yield a much greater pool of data from which to make inferences about the encoding of value by the midbrain dopaminergic system.

\subsection{A distributed code for dopaminergic signalling}\label{subsec:da_distributed_code}

It is a seductive metaphor to think of the brain as a kind of computer, and therefore that neurons are wires carrying a signal from one location in the brain to another. However, neuronal cells within the same information pathway are known to have considerable heterogeneity in their responses (\cite{martinFunctionalHeterogeneityNeighboring2013}; \cite{clancyStructureSingleWhisker2015}) and also perform computations at the individual cell level (\cite{brunelSingleNeuronDynamics2014}). Dopaminergic neurons in the midbrain that encode reward have also been shown to have such heterogeneity (\cite{ferrari-tonioloCodingBasicComponents2025}) and encode different aspects of probabilistic reinforcement learning, with individual neurons varying in their degree of optimism about future rewards (\cite{dabneyDistributionalCodeValue2020}, \cite{Lowet_Zheng_Meng_Matias_Drugowitsch_Uchida_2025}).

This heterogeneity raises important methodological considerations for studying dopaminergic signalling. If dopamine neurons implement a distributed code in which individual cells vary in their response properties and computational roles, then standard approaches that average responses across trials may inadvertently flatten this heterogeneity. Rather than every neuron performing the canonical reward prediction error (RPE) computation identically on each trial, individual neurons may contribute distinct components to the overall population signal. Trial-by-trial analysis of single neurons, rather than trial-averaged population responses, may therefore be necessary to uncover the true computational diversity within the dopaminergic system and understand how heterogeneous individual signals combine to produce reliable population-level representations of reward value.

\subsection{What does the dopamine signal actually compute?}\label{subsec:what_does_da_compute}

The work discussed in Section~\ref{subsec:da_distributed_code} concerns how dopamine encodes information, either as a single scalar signal or as a richer, distributed code across a population of neurons. It does not, however, settle \textit{what} that signal is actually encoding, which is a question we return to here. Classical reward prediction error theory predicts that once an outcome is fully and reliably predicted, the error signal should extinguish, since there is nothing left to be surprised about. Yet in Section~\ref{subsec:stereotyped_rpe}, we reported that dopaminergic responses reflecting the monkeys' bids, their trial profit, and their final liquid payment persisted even though the animals theoretically had perfect information about the forthcoming reward from the moment the computer's bid was revealed. We offered a partial explanation at the time -- the substantial delay between trial outcome and payment, together with the heterogeneity of individual BDM trials, which makes the calculations involved difficult. Similar persistent responses to well-predicted outcomes have been reported widely elsewhere in the dopaminergic literature, most notably as dopamine ``ramps'' that build steadily toward a reward whose timing and magnitude are already known to the animal (\cite{howeProlongedDopamineSignalling2013}; \cite{hamidMesolimbicDopamineSignals2016}), and as dopamine transients that are necessary and sufficient for learning that goes beyond simple, fully-extinguishable associations (\cite{sharpeDopamineTransientsAre2017}). This persistence appears to be a general property of the dopaminergic system, not a peculiarity of our own task or dataset -- a conclusion increasingly reflected in recent reviews of the field (\cite{kahntCuriousCaseDopaminergic2025}).

The central proposal of this thesis is that the resolution of this potential tension lies in recognising that dopamine tracks subjective utility, not the objective predictability of physical reward magnitude. That is, the dopaminergic cells encode the value to a monkey of some juice to be received, not just the volume. Having ``perfect information'' about the forthcoming payment, as our monkeys did, is not the same as there being no value-relevant variance left to signal: the animal's own subjective valuation, indexed by its bid, continued to fluctuate from trial to trial even though the physical reward on offer did not. This is a natural extension of the utility-prediction error account developed in Chapter 1 (\cite{staufferDopamineRewardPrediction2014}; \cite{lakDopaminePredictionError2014}), and is independently reinforced by a concurrent analysis from our own laboratory, which found that the dopamine population signal integrates both utility and subjectively weighted probability into a single measure of subjective value (\cite{ferrari-tonioloCodingBasicComponents2025}). It is also consistent with more recent treatments of prediction error itself. Rather than abandoning the concept, it has instead been argued that dopamine's error signal should be understood as operating over a much richer representation of state and value than the fixed, physical targets assumed by the earliest models (\cite{gershmanExplainingDopaminePrediction2024}). One concrete example of this comes from a model in which a cue's timing is learned rather than fixed in advance, and which reproduces dopamine responses that otherwise look like violations of a strict prediction-error rule (\cite{coneLearningExpressReward2024}). It is possible that the persistence we reported in Section~\ref{subsec:stereotyped_rpe} was never a puzzle to be explained away by trial heterogeneity: a response that tracks subjective value should not be expected to extinguish just because the objective payment became predictable.

An alternative, and increasingly prominent, explanation has been put forward by the Namboodiri laboratory and colleagues, who argue that dopamine is better understood as a retrospective, causal-learning signal than as a moment-to-moment prediction error over value (\cite{jeongMesolimbicDopamineRelease2022}; \cite{namboodiriWhyDopamineCausal2024}). On this account, dopamine tracks how reliably one event actually causes or predicts another, calculated by looking backwards from a meaningful outcome to work out what was responsible for it, rather than tracking how much value an anticipated outcome carries. This is a genuinely different kind of claim from the utility account above: it is a proposal about which events get credited as causes, not about how much any given event is worth. The strongest evidence for this proposal comes from a study in which dopamine responses to a cue-reward pair were selectively reduced when the statistical contingency between the cue and the reward was degraded, even though the timing between them, and therefore the temporal-difference prediction, was left completely unchanged (\cite{garrMesostriatalDopamineSensitive2024}). A purely value-based account struggles to explain this finding, which a causal-contingency framework predicts directly.

A further, longstanding distinction in the dopamine literature is between ``wanting'' and ``liking'' -- the motivational pull toward a reward and the hedonic pleasure taken in consuming it -- which Berridge and colleagues have argued are dissociable both behaviourally and neurally, with mesolimbic dopamine implicated specifically in the former (\cite{robinsonIncentiveSensitizationTheoryAddiction2025}). On this account, dopamine signals the incentive salience of a cue or outcome rather than a prediction error over its value as such (\cite{berridgePredictionErrorIncentive2012}), and this incentive can separate from, and even run contrary to, an animal's own learned predictions about how good an outcome will be (\cite{berridgeSeparatingDesirePrediction2023}). Consistent with a functional split of this kind, medium spiny neurons within the nucleus accumbens have recently been shown to carry distinct signals at the single-cell level, with one population responding to the salience of a stimulus regardless of its predicted value and another tracking prediction error more conventionally (\cite{zachryD1D2Medium2024}). Our own bids are naturally read as an economic report of value rather than of momentary wanting, and our task was not designed to dissociate the two; we discuss the fuller history and evidence for this distinction, including its computational reconciliation with prediction-error accounts, in \textbf{Appendix B}.

Our task was designed to elicit trial-by-trial variation in subjective value, not to manipulate the causal structure of the task itself. The link between the fractal cue, the bid, and the eventual reward was constant and unambiguous throughout every trial the monkeys completed. Our results therefore speak to how much value dopamine assigns to an anticipated outcome, but they are largely silent on the separate question of how reliably that outcome's cause is being tracked, which is the crux of the causal-learning account. The two accounts are not necessarily in competition: a system that assigns credit to causes could plausibly weight that credit by subjective value, and testing such a hybrid account directly would be a natural extension of this work. Other recent proposals extend this landscape in different directions, without directly competing with our own result. These include a dopamine signal proposed to track actions rather than value in a different pathway and species (\cite{greenstreetDopaminergicActionPrediction2025}), one framed around information gain rather than prediction error (\cite{beckPredictionErrorCorrelates2026}), and one in which dopamine's error signal is computed against the animal's own uncertain beliefs about the task (\cite{starkweatherDopamineRewardPrediction2017}). Taken together, our findings most directly extend the utility-coding tradition in which this thesis sits; how that tradition ultimately relates to causal and other accounts of dopaminergic function is a question this work raises rather than resolves.

\clearpage
\section{Concluding Remarks}

For a variety of reasons, this thesis took a long time in writing. Over the course of the years procrastinating, there have therefore been many times when different people have asked some version of the question 'Why study how dopamine contributes to economic decision-making?'

Clinically, we touch on disorders of the dopaminergic system in the first chapter of this thesis. In particular, in human patients, Parkinson's syndrome can lead to potentially catastrophic decision-making (\cite{gilmourImpairedValuebasedDecisionmaking2024}; \cite{denbrokApathyParkinsonsDisease2015}). Recently, there has also been an intense interest in how reward-seeking behaviour could be altered by pharmaceutical agents to reduce obesity in humans (\cite{wangGLP1ReceptorAgonists2023}). Unsurprisingly, there is therefore interest in how these drugs might modulate dopaminergic signalling in the human midbrain (\cite{zhuHedonicEatingControlled2025}; \cite{krupaCurbingAppetitesRestoring2025}).

However, we believe that studying how the brain manages to make economic decisions is interesting and vital beyond the scope of its clinical application. We are each second making value judgments which can determine anything from whether we have a pleasurable lunch, make a wise investment, or even live a fulfilling life.

Within the results of this thesis, we present evidence that it is possible to observe the dopaminergic representation of the value of a fractal indicating a magnitude of juice to a monkey, before that animal indicates its value to us via its bid. It is remarkable that we understand the brain to the extent that we can begin to uncover the shadow that lies between the motion and the act. However, it is not yet enough. Returning to the question at the beginning of this section: Why should we study how the brain produces economic decisions? Because we have these things in our head; they make us do things \textit{and we don't know why}. 